%!TEX root = ../../csuthesis_main.tex
\chapter{绪论}

\section{研究背景及意义}

物理信息神经网络（Physics-Informed Neural Networks, PINN）是一种将深度学习与先验物理知识相结合的机器学习框架，旨在将神经网络应用于传统数值计算领域，通过其强大的拟合能力和泛化能力解决与偏微分方程（Partial Differential Equations, PDE）相关的各种问题，包括方程求解、参数反演、模型发现、控制与优化等。

相比于传统的离散数值方法，如有限差分法 \cite{grossmann2007numerical}、有限元法 \cite{bathe2006finite}等，PINN求解偏微分方程 \cite{raissi2019physics}不再需要对求解域进行网格离散，只需要在求解区域内随机选取一定数量的数据点进行训练，简化了问题的预处理工作，尤其适用于具有复杂几何形态、动态边界或难以进行网格化的物理系统。另外，传统数值方法求解偏微分方程时依赖严格边界信息，而PINN则通过将物理定律融入损失函数，形成“软约束”，在仅有求解域内部数据或部分边界信息的情况下，仍能有效逼近真实解。除此之外，利用PINN对不同类型的偏微分方程进行求解时，不需要特定的设计，只需要利用最小化损失函数来迫使神经网络满足偏微分方程即可。上述特性使得PINN在流体力学 \cite{KQDX202308007,jin2021nsfnets}、材料科学 \cite{JSJG202304012,YDSG202305007,fang2019deep}和生物医学 \cite{LXXB202305014,kissas2020machine}等多个领域展现出巨大的应用潜力。

尽管PINN在理论层面和众多应用案例中显示出了巨大的潜力，但在实际应用中，PINN模型的训练过程往往面临着收敛速度慢、计算效率低的问题，尤其是在处理高维度、大规模或非均匀采样的复杂物理系统，如求解复杂流场时 \cite{wang2022and,fuks2020limitations}。传统PINN训练方法采用梯度下降优化算法，目标是最小化包含数据拟合误差与物理定律约束的复合损失函数。然而，在整个定义域上执行密集采样以满足物理定律及边界条件，使得PINN需要频繁进行前向传播和反向传播运算，加剧了计算复杂性和耗时问题。此外，由于物理规律通常以非凸形式体现在损失函数中，可能导致梯度消失或爆炸，严重影响模型训练的效率和最终解的精度。对于涉及多个子系统相互作用、强耦合关系或非线性动态特性的物理问题，PINN模型的收敛速度会进一步降低。因此，寻求有效的方法来加速PINN模型的训练过程，并提高其解决复杂物理问题的能力，是当前研究的重要方向。

\section{研究现状}

通常PINN求解的是一个高度非线性、非凸的最优化问题，而且涉及到对输入变量的高阶自动微分。大量的学者从不同维度对PINN模型进行了深入优化与改进，旨在提升其在解决实际物理问题时的精度与效率。这些方法可归纳为六个主要类别，本节将分别概述。

\subsection{域分解}

在PDE的数值求解过程中，单个神经网络能表达的容量是有限的，当面对具有复杂区域特征的问题时，需要在更复杂的区域上对解进行表达，一种直观的策略是通过增大网络规模，即增加深度与宽度，从而提升其表达复杂模式的能力。然而，在实际应用中，特别是对于PDE数值求解常用的多层感知机（Multi-Layer Perceptron, MLP）型神经网络，此类操作往往伴随着梯度爆炸或梯度消失等优化难题，导致求解过程变得异常困难，甚至影响最终解的精度与稳定性。因此，结合传统PDE数值求解方法中的域分解技术，通过构建多个专门针对不同区域进行学习的神经网络，能够有效降低单个网络的复杂度与学习难度，同时提高整体求解性能。

在求解时间依赖问题时，根据时域进行划分是一种常见的域分解技术。Meng等 \cite{meng2020ppinn}借鉴并行算法的思想，将一个长时间段的PDE问题分解成多个相互独立的短时段子问题，并通过快速且高效的粗粒度求解器进行全局监督与求解。空间域上的区域分割则是将求解域划分为多个子空间进行独立分析。例如，Moseley等 \cite{moseley2023finite}借鉴经典有限元方法，通过将输入空间细分为多个重叠小区域，并在每个区域内使用小型神经网络学习一组紧支撑的有限基函数之和来表示解。面对更加复杂的多尺度、多物理场问题，还可以采用更为综合与灵活的域分解技术。Jagtap等 \cite{jagtap2020extended}提出的扩展PINN（eXtended PINN, XPINN）允许在空间和时间两个维度上进行任意形式的分解。在XPINN框架下，每个子域采用一个独立的神经网络模型，并针对该子域的特点选择超参数，如网络深度、宽度、残差点的数量及其位置、激活函数以及优化算法等。对于解复杂度较高的子域，可以使用深层网络捕捉复杂的时空模式，而对于解相对简单平滑的子域则可采用浅层网络。这种灵活的域分解技术使得XPINN具有更强的表示能力和并行计算潜力，进一步提高了求解复杂PDE的效率和精度。

\subsection{采样方法}

在PINN中，偏微分方程损失通常用分散的时空点（称为残差点）来进行评估，因此，这些残差点的位置和分布对 PINN 的性能起着非常重要的作用。在训练过程中，残差点通常随机分布在域中。尽管该分布方式对于大多数情况都有效，但对于某些呈现陡峭梯度解的偏微分方程可能效率不高。以Burgers方程为例 \cite{lu2021deepxde}，应该在锐锋附近放置更多的点以更好地捕获不连续性。因此有学者对残差点采样策略进行了详细研究。

Wu等 \cite{wu2023comprehensive}系统地探讨了两类PINN中的关键采样方法：非自适应均匀采样和自适应非均匀采样,并考虑了六种不同的均匀采样方法：等间距网格、均匀随机采样、拉丁超立方采样、Halton 序列、Hammersley序列、Sobol 序列以及一种重采样策略（Random-R）。同时，为了提高PINN的采样效率和准确率，Wu等又提出了两种新的基于残差的自适应采样方法：
\begin{itemize}
    \item 基于残差的自适应分布（Residual-based Adaptive Distribution, RAD）；
    \item 结合分布的基于残差的自适应细化（Residual-based Adaptive Refinement with Distribution, RAR-D）。
\end{itemize}
然后，通过大量的数值实验对比分析了不同采样策略在四个正向问题和两个逆向问题上的性能表现。结果表明，作者所提出的RAD方法和RAR-D方法在所有测试问题上始终表现最优。这些结果为选择PINN中合适的采样方法提供了实用的指导依据。

由于RAR及其变体方法通过蒙特卡洛积分方法密集采样来计算PDE残差分布，产生了较大的计算负担，Daw等 \cite{daw2022mitigating}提出一种新的采样策略——保留-重采样-释放方法（Retain-Resample-Release, R3）。该方法通过保留高残差区域的采样点、重新采样并更新点集，可以在高残差区域中增量累积搭配点，而几乎没有计算开销。

针对时间依赖问题，Guo等 \cite{guo2022novel}指出仅依赖PDE残差进行自适应采样并不充分，采样过程还应遵循时间因果性原则。为此，他们将时间因果性引入自适应采样策略中，设计出新的自适应因果采样方法（Adaptive Causal Sampling Method, ACSM）。        

\subsection{损失权重}

在PINN的训练过程中，通常使用基于梯度的 Adam（Adaptive Moment Estimation） \cite{kingma2014adam}和 L-BFGS（Limited-memory Broyden-Fletcher-Goldfarb-Shanno） \cite{liu1989limited}优化器，但由于损失函数的各项梯度不平衡，导致训练困难和模型性能下降。动态权重调整策略可以有效缓解上述问题，并提升模型对PDE解的精确拟合和整体预测性能。

一些学者对PINN的训练过程进行了分析，并使用分析得到的统计数据对权重进行系统地调整。Wang等 \cite{wang2021understanding}通过研究PINN训练过程中的梯度变化，揭示了残差损失梯度和边界损失梯度两个损失项之间不平衡的机制，并利用该梯度统计信息自动调整两者的权重，从而提出了一种学习率退火算法。尽管该算法能够有效改善损失函数中不同项之间的平衡问题，但缺乏对PINN训练的理论指导，因此Wang等 \cite{wang2022and}利用神经切线核（Neural Tangent Kernel, NTK）理论研究了PINN训练过程中的梯度传播问题，并且利用NTK的特征值来分配损失项的权重，从而提出了一种新的梯度下降算法，该方法能够有效实现各损失分量收敛的均衡，提高PINN的训练性能。

在PINN的训练过程中，损失权重应当最大化以强调不同损失项的重要性，而神经网络权重则被调整以最小化损失函数。因此，PINN的自适应加权训练可以被转化为最小最大值问题。Liu等 \cite{liu2021dual}提出了一种具有最小最大结构的新型物理约束神经网络（Physics-Constrained Neural Networks with Minimax Architecture, PCNN-MM），核心思想是利用双二聚体方法来搜索非凸-非凹目标函数的高阶鞍点。Levi等 \cite{mcclenny2020self}采用可训练的损失权重来生成自适应掩码函数，类似于计算机视觉中使用的注意力机制，这些自适应权重通过反向传播和梯度上升实现最大化，同时神经网络网络权重进行梯度下降以实现最小化。Xiang等 \cite{xiang2022self}利用高斯概率模型定义各个损失项，并引入噪声参数集合作为损失项的权重系数。在每个训练迭代中，通过最大似然估计来更新这些噪声参数，从而实现了损失函数的自动调整和平衡。

\subsection{优化算法}

根据1.2.3节中的介绍，大部分学者通过损失重加权来估计损失的相对重要性，也有学者将该思路直接应用于参数优化环节。Yao等 \cite{yao2023multiadam}提出了MultiAdam优化器，利用梯度动量在参数级别上平衡PDE损失和边界损失，以应对两者间的不平衡问题。通过这种方式，实现了优化过程中的尺度不变，从而增强模型对不同尺度问题的适应性。另一种改进的思路则是直接采用不同的优化算法，Jnini等 \cite{jnini2024gauss}使用函数空间中的高斯-牛顿方法求解Navier-Stokes方程，达到了极高的求解精度。

\subsection{新损失函数}

在损失函数方面的改进除了调整损失权重外，还有学者对损失函数本身的组成进行了改进。Yu等 \cite{yu2022gradient}提出了梯度增强物理信息神经网络（Gradient-enhanced PINN, gPINN）。在PINN的基础上，gPINN不仅关注PDE残差本身，还关注其梯度，通过增加一个额外的损失项来约束PDE残差的梯度，确保在PDE残差为零时，其梯度也为零。Kharazmi等 \cite{kharazmi2021hp}则提出了一种名为“hp-变分物理信息神经网络”的通用框架，在构建损失函数时，遵循加权残差法，使用变分损失函数以指导神经网络的学习过程。

\subsection{新网络结构}

大部分的综述 \cite{lu2021deepxde, 查文舒2022基于神经网络的偏微分方程求解方法研究综述,cuomo2022scientific}将PINN归类到Deep Learning中，也就是需要多层的神经网络，多层的网络是造成PINN本身非线性非凸从而难以求解的原因之一。为了更快速地求解PINN，又有学者 \cite{dwivedi2020physics}将目光放回到浅层网络中，也就是结合了极限学习机（Extreme Learning Machine, ELM）的物理信息极限学习机（Physics-Informed ELM, PIELM），仅包含一个隐藏层，而且隐藏层的输入参数进行随机初始化后被固定不再需要更新。Dwivedi等 \cite{dwivedi2020physics}就提出了一种分布式的DPIELM算法，通过在不同区域使用不同的表示并施加连续性和可微性约束，提高了对复杂函数的表征能力和解决时变问题的效果。Bai等 \cite{bai2023physics}深入研究了PINN在训练过程中呈现出局部逼近器的特性，并基于这一观察设计了在整个训练过程中能够保持局部逼近性质，且仅包括一个隐藏层和采用径向基函数作为“激活”功能的单隐藏层神经网络。

\subsection{融合数值微分}

传统的PINN通过自动微分（Automatic Differentiation，AD）来计算损失函数中的导数项，并结合偏微分方程约束训练过程以确保输出结果符合物理定律。然而，仅依赖AD的PINN在采样点不足的情况下容易导致优化趋向于非物理解，即需要大量的离散点才能实现高精度。

为了解决这一问题并提高训练效率与准确性，Chiu等 \cite{chiu2022can}通过结合数值微分（Numerical Differentiation, ND），设计了一种耦合自动-数值微分（Coupled-Automatic–Numerical Differentiation）框架，构建了CAN-PINN模型。CAN-PINN利用局部支持点上的值以及AD计算得到的导数项，在稀疏采样情况下能有效关联临近的离散点，从而增强训练效率，并且相比基于ND的PINN能够将精度提升1到2个数量级。Sharma等 \cite{sharma2022accelerated}则使用无网格径向基函数有限差分法对PINN损失函数中的高阶空间导数进行高效且高精度的数值离散化替代，显著降低了自动微分过程中重复计算高阶导数所带来高昂计算成本。Fang等 \cite{fang2021high}结合卷积神经网络和有限体积法的思想，通过局部拟合方法对微分算子进行近似，可以避免传统PINN中因神经网络预测不佳而导致的问题。

\section{研究内容}

基于上述研究现状，本文具体研究内容如下。

第一章是绪论，主要介绍关于物理信息神经网络的研究背景，阐述PINN相对于传统数值离散方法的优势，分析当前PINN模型训练过程遇到的问题。在研究现状一节，从不同方向总结PINN的优化与改进方法，包括域分解技术、自适应采样策略、动态损失权重调整、优化算法的改进、新型损失函数的设计、新网络结构的开发以及数值微分与自动微分相结合的策略。研究内容一节对本文的研究问题和方法进行总结和陈述。

第二章深入探讨物理信息神经网络模型，从深度神经网络的基础与发展历程出发，阐述其基本原理、初始化方法与优化算法。PINN通过将物理定律融入神经网络的损失函数，能够在数据驱动的基础上结合物理先验知识求解偏微分方程。与传统机器学习相比，PINN在融合物理信息、降低数据依赖性、增强泛化能力和处理特定物理问题方面展现出明显优势。本章还详细阐述PINN在求解Navier-Stokes方程的应用案例，包括模型设计、训练细节与结果分析，证明其在复杂流体动力学仿真中的有效性。最后，针对PINN存在的局限性，提出模型改进优化的方向：采用自适应采样策略、动态调整损失权重和改进网络结构。

第三章主要讨论自适应采样方法。首先回顾传统的固定采样与随机重采样方法，指出其在复杂PDE问题中的局限性。接着，深入分析基于残差的自适应采样策略：基于残差的自适应细化和自适应分布，前者在残差较大的区域动态增加采样点，后者通过PDE残差分布来调整重采样的概率密度，两者均旨在优化采样点布局以更好地拟合PDE解。随后，本章又引出“保留-重采样-释放”方法，旨在克服RAD方法中计算成本高和参数调整难等问题。通过对比不同采样策略在迭代过程中的$L^2$相对误差，验证R3算法在模型收敛速度和最终精度上的显著优势。

第四章主要讨论基于残差的自适应物理信息神经网络方法，旨在克服传统PINN在求解复杂偏微分方程时的收敛性和精确性挑战。首先介绍一种自适应权重策略，通过可训练权重自动平衡不同损失项。接着阐述残差注意力方法，根据累积残差动态调整残差点权重。此外，本章还引入一种基于Transformer的改进多层感知机架构，通过结合残差连接与自注意力机制增强神经网络对输入特征的交互捕捉和深层表示能力，提升模型性能。实验部分通过比较不同策略在收敛速度和预测精度上的影响，验证所提方法的有效性。

第五章对全文进行总结，提出目前物理信息神经网络优化加速中亟待解决的问题，并对未来进一步的研究进行了展望。
